\scanpage{152}
(NB $K_r=\bigcap_i\operatorname{Ker}(\sigma_i,\overline\sigma_i:
K\rightrightarrows\overline{\mathbb Q})$ est stable par tt changement de base
$\mathbb Q\longrightarrow\Lambda$). Donc
\[
\operatorname{Ker}\psi_\Sigma\cap\bigcap_{1\leq i\leq d'}\operatorname{Ker}\psi_{\Sigma_i}
=1\qquad\text{dans }S_K^0
\]

\uline{Question} Regardons les hom. $\psi_\Sigma:S_K^0\longrightarrow T_E$, avec $E$
\uline{quadratique}. Y en a-t-il assez, du moins en filtrant sur $K$ croissants?
Or un $E$ a, relativt à une \deleted{\unsure{E}} extension $K$ de $\mathbb Q$,
\deleted{deux} $0$ structure CM si $E$ ne pas plonger dans $K$, et deux exactement
dans le cas contraire \inserted{(correspond. aux plongts)}
\inserted{[L'hom $\psi_K$ correspondant est \deleted{\operatorname{Tr}_{K/E}} $N_{K/E}$.}
Donc pour $K$ fixé, les hom $\psi_\Sigma:S_K^0-T_E$ correspondent par paires conjuguées
aux plongts $E\longrightarrow K$; donc pour voir s'il y en a assez, il faut voir s'il y a
``assez'' de sous-corps quadratiques de $K$.
Si $K$ \inserted{(supposé Weilien)} n'est pas \inserted{(dés. $L$ ce ss-corps)}
engendré par ses sous-corps quadratiques, i.e. si \deleted{le groupe de} \inserted{$K$ n'est}
pas galoisien de groupe de Galois $(\mathbb Z/2\mathbb Z)^N$, alors le noyau de
$N_{K/L}:S_K^0-S_L^0$ est un tore de dim $>0$ qui est dans le noyau de tous les
$\psi_K$ envisagés.
\margin{151}
